\section{A moving surface with separate virgin, char and gas inventories}
The candidate transports the masses of virgin phenolic $M_v$, char $M_c$, and
their combined sensible enthalpy $H_s$ in each cylindrical cell. A cell has
volume $(r_o^2-r_i^2)/2$ per radian per metre of axial length. Conservative
overlap integration in $r^2/2$ returns both the retained and swept inventories
as the face advances. In the equations below, $M$ and $H$ may refer either to
a physical cell or to these quantities per radian per metre, provided every
source uses the same convention.

\subsection{Bulk reaction and reconstruction}
We use a common reference $T_*=22\,{}^\circ$C for virgin solid, char and gas.
The inherited table defines piecewise-linear specific heat capacities. Let
$h_i(T)=\int_{T_*}^T c_{p,i}(u)\,du$ for $i=v,c$ and
$h_g(T)=1600(T-T_*)$ J/kg when $T$ is in degrees Celsius. The candidate
holds the first heat capacity constant between 0 and 22 degrees Celsius and
rejects states above the final 2476.85 degree table entry or below zero.
Positive phase masses give a strictly increasing enthalpy relation
$H_s=M_vh_v(T)+M_ch_c(T)$. Inverting the quadratic integral within its
table interval recovers a single temperature after remapping. An exhausted
cell carries zero mass and enthalpy and has no temperature to reconstruct.

Let $y=600/1250=0.48$ denote the assumed char yield per kilogram of virgin
material consumed. With $\bar T=(T^n+T^{n+1})/2$ and $k=333\exp[-64081/
(R(\bar T+273.15))]$ in reciprocal seconds, the update uses
\[
 d=M_v^n(1-e^{-k\Delta t}),\quad
 M_v^{n+1}=M_v^n-d,\quad M_c^{n+1}=M_c^n+yd,\quad m_g=(1-y)d .
\]
Mass and enthalpy use the same $d$. The scalar temperature solve enforces
\[
 H_s^{n+1}-H_s^n=Q_{\rm cell}-L_p d-m_g h_g(\bar T).
\]
The chosen basis for $L_p=418$ kJ/kg is virgin mass consumed at $T_*$. Gas
leaves carrying its sensible enthalpy at generation. The phase inventories on
the left already include the change in stored solid sensible heat.

\subsection{Partially converted material at the wall}
Let $a$ and $b$ be the virgin and char mass densities in the material adjacent
to the gas. The current candidate assumes that residual virgin material swept
by the front completes pyrolysis there and its char reacts at that same face.
For speed $v\geq0$ the outgoing surface pyrolysis gas flux is $j_g=(1-y)av$
and the outgoing char-product flux is $j_c=(b+ya)v$, in kg/(m$^2$ s).
Their sum equals the rate $(a+b)v$ at which the surface sweeps up solid.
Writing $\alpha=b/(b+ya)$, we impose $j_c=\alpha j_{\rm law}(T_s,p)$ with
the inherited kinetic and mass-transfer law. A virgin wall can heat and
pyrolyse volumetrically before surface recession begins.

For a first-cell temperature $T_1$, define the sensible power delivered to
the boundary by swept inventory as $q_{\rm swept}=v[ah_v(T_1)+bh_c(T_1)]$.
Per unit hot-face area the wall balance is
\begin{align*}
 q_{\rm conv}+q_{\rm rad}+q_{\rm swept}
 &=q_{\rm cond}+j_gh_g(T_s)+j_ch_c(T_s)+avL_p+j_cL_c,\\
 q_{\rm cond}
 &=\frac{\displaystyle\int_{T_1}^{T_s}[(1-\alpha)k_v(T)+\alpha k_c(T)]\,dT}
 {r_s\log(r_1/r_s)} .
\end{align*}
Here $q_{\rm cond}>0$ flows into the retained solid, $r_s$ is the moving wall
radius and $r_1$ the adjacent cell centre. The conduction expression integrates
the cylindrical resistance at fixed local composition. The film coefficient
uses total outward mass flux, including bulk gas arriving at the surface and
both $j_g$ and $j_c$. A bracketed scalar solve closes film heating, optional
radiation, conduction and recession at the same $T_s$. Bulk gas has already
exported its sensible heat in the bulk equation. The chosen proxy $h_c(T_s)$
for the char product and excess cost $L_c=35$ MJ/kg together specify a
provisional removal energy. Data for reaction products, recession and pyrolysis
are needed to determine whether these assumptions hold for the Iris liner.

\subsection{Verification to date}
Tests compare the enthalpy primitive against independent quadrature and the
isothermal reaction against the analytic solution. In a four-second manufactured
heated cell, doubling the time steps from 8 to 16, 32 and 64 reduces temperature
error against an independent ODE integration from 0.151 to 0.0376, 0.00940
and 0.00235 K. A nonuniform reacting shell with prescribed recession closes
the integrated mass and energy ledgers through final-cell removal to relative
errors below $2\times10^{-15}$. Separate manufactured tests recover imposed
face temperatures for a virgin conducting wall, pure reacting char, and a
mixed wall with blowing and radiation. The face equation residual is below
$3\times10^{-9}$ W/m$^2$ for those fixtures. These checks exercise individual
material and boundary balances. The coupled transient wall, native
implementation and event-time comparison remain open.
